% 書前：考試出題分析與學習路徑（不編章號）。
% 數字由私有資料產生（aki-judge-problems/dsa-book/exam-analysis/gen_guide.py），不要手改圖表裡的數字。
\chapter*{考試出題分析與學習路徑}
\addcontentsline{toc}{chapter}{考試出題分析與學習路徑}
\markboth{考試出題分析與學習路徑}{}

這一章不教內容，只回答兩個問題：研究所考試到底考哪些東西，以及你該照什麼順序讀這本書。

\section*{資料從哪裡來}
\addcontentsline{toc}{section}{資料從哪裡來}

我們收集了 10 所學校公開的研究所入學考試考卷，學年度從 77 到 115，有考資料結構或演算法的共 324 份，3543 題。 考卷分兩類：

\begin{itemize}
\item \textbf{資工所}（202 份）：資工、資科、資訊聯招的「資料結構與演算法」主科，以及資料結構和其他科目印在一起的合卷、含資料結構的計算機概論。
\item \textbf{他系}（122 份）：電機、電信、工科、資管、統計等系所單考「資料結構」的考卷。
\end{itemize}

每一題都對到本書最接近的一節（一題考兩個主題就兩節都算）。 合卷裡作業系統、計算機組織這類題目不算。 後面用兩個數字：

\begin{itemize}
\item \textbf{出現率}：有考到這一章的考卷，佔全部考卷的比例。 一份考卷考了三題堆積，堆積也只算一次。
\item \textbf{配分比例}：這一章的題目，佔所有有標配分題目總分的比例。 一題同時考兩章，分數各算一半。
\end{itemize}

出現率回答「會不會考」，配分比例回答「考多重」。 兩個都高的章，是最划算的。

\section*{哪些章最常考}
\addcontentsline{toc}{section}{哪些章最常考}

「各章出現率」圖是每一章的出現率，紅色是資工所，灰色是他系。

\begin{figure}[tbp]
\centering
\textbf{各章出現率}\par\smallskip
\begin{tikzpicture}
\begin{axis}[
  xbar, width=0.82\linewidth, height=17cm, bar width=4pt,
  y dir=reverse, ytick=data, symbolic y coords={{1 分析與漸近符號},{2 遞迴式},{3 平攤分析},{4 陣列與矩陣},{5 鏈結串列},{6 堆疊與佇列},{7 樹與二元樹},{8 堆積},{9 二元搜尋樹},{10 AVL 與紅黑樹},{11 多路搜尋樹與 trie},{12 雜湊},{13 互斥集合},{14 排序},{15 下界與線性排序},{16 選擇問題},{17 分治},{18 動態規劃},{19 貪心},{20 回溯與分支限界},{21 BFS 與 DFS},{22 拓樸排序與 AOE},{23 連通性},{24 最小生成樹},{25 最短路徑},{26 網路流},{27 字串比對},{28 NP 完備},{29 近似演算法}},
  yticklabel style={font=\small}, xmin=0, xmax=100,
  xlabel={出現在幾 \% 的考卷裡}, xmajorgrids, grid style={black!12},
  enlarge y limits=0.02, legend style={at={(0.98,0.02)},anchor=south east,font=\small},
  legend cell align=left,
]
\addplot[fill=examStroke!80, draw=none] coordinates {(35,{1 分析與漸近符號}) (53,{2 遞迴式}) (6,{3 平攤分析}) (14,{4 陣列與矩陣}) (27,{5 鏈結串列}) (48,{6 堆疊與佇列}) (57,{7 樹與二元樹}) (47,{8 堆積}) (37,{9 二元搜尋樹}) (42,{10 AVL 與紅黑樹}) (24,{11 多路搜尋樹與 trie}) (42,{12 雜湊}) (7,{13 互斥集合}) (60,{14 排序}) (30,{15 下界與線性排序}) (7,{16 選擇問題}) (26,{17 分治}) (50,{18 動態規劃}) (29,{19 貪心}) (7,{20 回溯與分支限界}) (48,{21 BFS 與 DFS}) (16,{22 拓樸排序與 AOE}) (12,{23 連通性}) (40,{24 最小生成樹}) (33,{25 最短路徑}) (10,{26 網路流}) (8,{27 字串比對}) (32,{28 NP 完備}) (5,{29 近似演算法})};
\addplot[fill=black!35, draw=none] coordinates {(52,{1 分析與漸近符號}) (23,{2 遞迴式}) (7,{3 平攤分析}) (33,{4 陣列與矩陣}) (57,{5 鏈結串列}) (75,{6 堆疊與佇列}) (80,{7 樹與二元樹}) (76,{8 堆積}) (54,{9 二元搜尋樹}) (47,{10 AVL 與紅黑樹}) (32,{11 多路搜尋樹與 trie}) (57,{12 雜湊}) (12,{13 互斥集合}) (71,{14 排序}) (33,{15 下界與線性排序}) (2,{16 選擇問題}) (19,{17 分治}) (25,{18 動態規劃}) (15,{19 貪心}) (8,{20 回溯與分支限界}) (77,{21 BFS 與 DFS}) (21,{22 拓樸排序與 AOE}) (30,{23 連通性}) (44,{24 最小生成樹}) (34,{25 最短路徑}) (4,{26 網路流}) (7,{27 字串比對}) (11,{28 NP 完備}) (0,{29 近似演算法})};
\legend{資工所（202 份）, 他系資料結構（122 份）}
\end{axis}
\end{tikzpicture}
\end{figure}

幾件看得出來的事：

\begin{itemize}
\item 樹、排序、堆積、堆疊與佇列、BFS 與 DFS，兩類考卷都有將近或超過一半的考卷會考。 這幾章是地基，怎麼準備都跳不掉。
\item 遞迴式、動態規劃、貪心、NP 完備，資工所考得比他系多很多；他系則多考鏈結串列、陣列、連通性（割點、雙連通元件）。 他系的考卷偏「資料結構」，資工所的偏「演算法」。
\item 平攤分析、互斥集合、選擇問題、回溯、字串比對、近似演算法，兩類考卷都不到一成。
\end{itemize}

\section*{哪些章拿分最多}
\addcontentsline{toc}{section}{哪些章拿分最多}

「資工所近 10 年各章配分比例」圖是 106--115 學年資工所考卷裡，各章佔總配分的比例，由高到低排。

\begin{figure}[tbp]
\centering
\textbf{資工所近 10 年各章配分比例}\par\smallskip
\begin{tikzpicture}
\begin{axis}[
  xbar, width=0.82\linewidth, height=18cm, bar width=6pt,
  y dir=reverse, ytick=data, symbolic y coords={{18 動態規劃},{14 排序},{2 遞迴式},{7 樹與二元樹},{25 最短路徑},{8 堆積},{6 堆疊與佇列},{24 最小生成樹},{21 BFS 與 DFS},{12 雜湊},{28 NP 完備},{10 AVL 與紅黑樹},{9 二元搜尋樹},{5 鏈結串列},{19 貪心},{1 分析與漸近符號},{15 下界與線性排序},{17 分治},{11 多路搜尋樹與 trie},{26 網路流},{22 拓樸排序與 AOE},{4 陣列與矩陣},{3 平攤分析},{23 連通性},{29 近似演算法},{20 回溯與分支限界},{27 字串比對},{13 互斥集合},{16 選擇問題}},
  yticklabel style={font=\small}, xmin=0, xmax=10,
  xlabel={佔總配分的 \%}, xmajorgrids, grid style={black!12}, enlarge y limits=0.02,
  nodes near coords, nodes near coords style={font=\scriptsize}, point meta=x,
  every node near coord/.append style={/pgf/number format/.cd, fixed, fixed zerofill, precision=1},
]
\addplot[fill=examStroke!80, draw=none] coordinates {(8.1,{18 動態規劃}) (7.9,{14 排序}) (7.7,{2 遞迴式}) (7.2,{7 樹與二元樹}) (6.0,{25 最短路徑}) (5.9,{8 堆積}) (5.6,{6 堆疊與佇列}) (5.4,{24 最小生成樹}) (5.3,{21 BFS 與 DFS}) (4.5,{12 雜湊}) (4.1,{28 NP 完備}) (3.7,{10 AVL 與紅黑樹}) (3.5,{9 二元搜尋樹}) (3.3,{5 鏈結串列}) (3.0,{19 貪心}) (2.7,{1 分析與漸近符號}) (2.0,{15 下界與線性排序}) (1.8,{17 分治}) (1.8,{11 多路搜尋樹與 trie}) (1.6,{26 網路流}) (1.5,{22 拓樸排序與 AOE}) (1.4,{4 陣列與矩陣}) (1.3,{3 平攤分析}) (1.1,{23 連通性}) (1.1,{29 近似演算法}) (0.8,{20 回溯與分支限界}) (0.7,{27 字串比對}) (0.7,{13 互斥集合}) (0.4,{16 選擇問題})};
\end{axis}
\end{tikzpicture}
\end{figure}

動態規劃、排序、遞迴式、樹這四章合起來超過三成的分數。 最短路徑的出現率只排中間，配分卻排第五：它常出 10 分以上的大題（手算 Dijkstra、Floyd--Warshall，或設計一個變形）。 反過來，漸近符號的出現率不低，但多半是一題幾分的選擇題，配分比例不高。

\section*{各章分級}
\addcontentsline{toc}{section}{各章分級}

照資工所近 10 年的出現率分成四級：\textbf{必考}（$40\%$ 以上）、\textbf{常考}（$20\%$--$40\%$）、\textbf{偶爾}（$10\%$--$20\%$）、\textbf{少考}（不到 $10\%$）。 「更早」是 105 學年以前的資工所考卷。

{\small
\begin{longtable}{rlcrrrr}
\toprule
章 & 主題 & 分級 & \multicolumn{2}{c}{資工所出現率} & 資工所近 10 年 & 他系 \\
\cmidrule(lr){4-5}
 & & & 近 10 年 & 更早 & 配分比例 & 出現率 \\
\midrule
\endhead
1 & 分析與漸近符號 & 常考 & 33\% & 36\% & 2.7\% & 52\% \\
2 & 遞迴式 & 必考 & 50\% & 56\% & 7.7\% & 23\% \\
3 & 平攤分析 & 少考 & 10\% & 2\% & 1.3\% & 7\% \\
4 & 陣列與矩陣 & 偶爾 & 13\% & 15\% & 1.4\% & 33\% \\
5 & 鏈結串列 & 常考 & 29\% & 26\% & 3.3\% & 57\% \\
6 & 堆疊與佇列 & 必考 & 51\% & 44\% & 5.6\% & 75\% \\
7 & 樹與二元樹 & 必考 & 60\% & 55\% & 7.2\% & 80\% \\
8 & 堆積 & 必考 & 55\% & 38\% & 5.9\% & 76\% \\
9 & 二元搜尋樹 & 常考 & 39\% & 35\% & 3.5\% & 54\% \\
10 & AVL 與紅黑樹 & 必考 & 43\% & 40\% & 3.7\% & 47\% \\
11 & 多路搜尋樹與 trie & 偶爾 & 16\% & 32\% & 1.8\% & 32\% \\
12 & 雜湊 & 必考 & 47\% & 37\% & 4.5\% & 57\% \\
13 & 互斥集合 & 少考 & 6\% & 9\% & 0.7\% & 12\% \\
14 & 排序 & 必考 & 67\% & 53\% & 7.9\% & 71\% \\
15 & 下界與線性排序 & 常考 & 26\% & 35\% & 2.0\% & 33\% \\
16 & 選擇問題 & 少考 & 8\% & 7\% & 0.4\% & 2\% \\
17 & 分治 & 偶爾 & 16\% & 36\% & 1.8\% & 19\% \\
18 & 動態規劃 & 必考 & 52\% & 47\% & 8.1\% & 25\% \\
19 & 貪心 & 常考 & 29\% & 29\% & 3.0\% & 15\% \\
20 & 回溯與分支限界 & 少考 & 10\% & 5\% & 0.8\% & 8\% \\
21 & BFS 與 DFS & 必考 & 43\% & 54\% & 5.3\% & 77\% \\
22 & 拓樸排序與 AOE & 偶爾 & 18\% & 13\% & 1.5\% & 21\% \\
23 & 連通性 & 偶爾 & 14\% & 9\% & 1.1\% & 30\% \\
24 & 最小生成樹 & 必考 & 43\% & 37\% & 5.4\% & 44\% \\
25 & 最短路徑 & 常考 & 39\% & 27\% & 6.0\% & 34\% \\
26 & 網路流 & 偶爾 & 14\% & 5\% & 1.6\% & 4\% \\
27 & 字串比對 & 少考 & 10\% & 6\% & 0.7\% & 7\% \\
28 & NP 完備 & 常考 & 36\% & 27\% & 4.1\% & 11\% \\
29 & 近似演算法 & 少考 & 8\% & 3\% & 1.1\% & 0\% \\
\bottomrule
\end{longtable}}

\textbf{近 10 年變多的}：第~8~章堆積（38\%→55\%）、第~14~章排序（53\%→67\%）、第~25~章最短路徑（27\%→39\%）、第~12~章雜湊（37\%→47\%）、第~28~章NP 完備（27\%→36\%）、第~26~章網路流（5\%→14\%）、第~3~章平攤分析（2\%→10\%）、第~6~章堆疊與佇列（44\%→51\%）、第~24~章最小生成樹（37\%→43\%）、第~7~章樹與二元樹（55\%→60\%）、第~23~章連通性（9\%→14\%）。 \textbf{變少的}：第~17~章分治（36\%→16\%）、第~11~章多路搜尋樹與 trie（32\%→16\%）、第~21~章BFS 與 DFS（54\%→43\%）、第~15~章下界與線性排序（35\%→26\%）、第~2~章遞迴式（56\%→50\%）。 堆積、雜湊、最短路徑、NP 完備、網路流、平攤分析都在往上；分治（舊卷常考二分搜尋的變形和計算幾何）、多路搜尋樹（舊卷常考 B 樹的磁碟存取細節和索引檔）在往下。

\section*{題型}
\addcontentsline{toc}{section}{題型}

\begin{center}
\begin{tikzpicture}
\begin{axis}[
  xbar, width=0.75\linewidth, height=7.5cm, bar width=5pt,
  y dir=reverse, ytick=data, symbolic y coords={{選擇／是非},{手算追蹤},{計算／複雜度},{設計演算法},{簡答觀念},{讀寫程式},{證明},{填空}},
  yticklabel style={font=\small}, xmin=0, xmax=50,
  xlabel={佔題數的 \%}, xmajorgrids, grid style={black!12}, enlarge y limits=0.08,
  legend style={at={(0.98,0.02)},anchor=south east,font=\small}, legend cell align=left,
]
\addplot[fill=examStroke!80, draw=none] coordinates {(35,{選擇／是非}) (18,{手算追蹤}) (14,{計算／複雜度}) (11,{設計演算法}) (10,{簡答觀念}) (4,{讀寫程式}) (4,{證明}) (2,{填空})};
\addplot[fill=black!35, draw=none] coordinates {(43,{選擇／是非}) (21,{手算追蹤}) (8,{計算／複雜度}) (6,{設計演算法}) (12,{簡答觀念}) (7,{讀寫程式}) (2,{證明}) (1,{填空})};
\legend{資工所, 他系資料結構}
\end{axis}
\end{tikzpicture}
\end{center}

三分之一以上是選擇題，資工所近 10 年更將近一半，而且不少學校用複選、答錯倒扣。 複選題每個選項都是一個小判斷，觀念要很精確，「大概知道」會被倒扣吃掉。 第二多的是手算追蹤：畫出插入後的樹、跑一遍演算法、填出表格。 這兩類加起來超過一半，所以本書每章的觀念題以這兩種為主，附錄~B 專門講追蹤題怎麼寫才不會被扣分。

證明題只佔 $4\%$ 左右，但大多集中在三四所學校，而且一題常是 10 分以上。 要考那些學校，第~1~章的迴圈不變量和附錄~B 的證明寫法要練。

\section*{各校偏重}
\addcontentsline{toc}{section}{各校偏重}

下表是資工所考卷份數 14 份以上的學校，配分比例最高的三章（括號是佔該校總配分的比例）。 份數少的學校，數字只能參考。

{\small
\begin{center}
\begin{tabular}{lrcl}
\toprule
學校 & 份數 & 學年度 & 配分最高的三章 \\
\midrule
成大 & 39 & 77--115 & 遞迴式（9\%）、動態規劃（9\%）、最短路徑（6\%） \\
中山 & 37 & 87--115 & 遞迴式（15\%）、樹與二元樹（12\%）、排序（12\%） \\
臺大 & 32 & 84--115 & 動態規劃（10\%）、遞迴式（9\%）、堆積（6\%） \\
中央 & 23 & 93--115 & 動態規劃（10\%）、BFS 與 DFS（9\%）、最短路徑（8\%） \\
政大 & 20 & 96--115 & 遞迴式（9\%）、鏈結串列（7\%）、動態規劃（7\%） \\
北科大 & 18 & 106--115 & 堆疊與佇列（13\%）、鏈結串列（10\%）、動態規劃（10\%） \\
中正 & 15 & 108--115 & 排序（13\%）、樹與二元樹（12\%）、最短路徑（11\%） \\
臺科大 & 14 & 102--115 & 樹與二元樹（18\%）、分析與漸近符號（16\%）、二元搜尋樹（11\%） \\
\bottomrule
\end{tabular}
\end{center}}

各校的差別比想像中小：每一校的前幾名都在「必考」那一級裡。 差別在比重，例如有的學校遞迴式一章就佔一成五，有的學校偏資料結構的實作（堆疊、串列）。 報考前，把目標學校最近三到五年的考卷做一遍，比任何統計都準。

\section*{學習路徑}
\addcontentsline{toc}{section}{學習路徑}

照你的目標選一條。 每條路徑的章都照建議順序排；「略讀」是讀定義和結論、做觀念題，不細讀證明和 C 程式；沒列出來的章可以先跳過。

\subsection*{路徑一：資工所，完整準備（約 4--6 個月）}

適合：考資工、資科、資訊聯招，時間夠，想要每一章都穩。

\begin{enumerate}
\item 基礎：第~1、2~章（主定理、遞迴樹一定要熟），第~3~章略讀。
\item 線性結構與樹：第~5、6、7、8、9、10、12~章。 每章的手算追蹤都要自己畫一次。
\item 排序：第~14、15~章。
\item 演算法設計：第~18~章（配分最高，花最多時間）、第~19~章、第~17~章。
\item 圖論：第~21、24、25~章，再第~22、23~章。
\item 理論：第~28~章（歸約的方向一定要會判斷）。
\item 有餘力再讀：第~11、26、13、16、20、27、29~章。
\end{enumerate}

\subsection*{路徑二：考前衝刺（約 6--8 週）}

適合：時間不多，只想把最常考、最多分的拿到。 只讀「必考」級，再加兩章配分高的常考章。

\begin{enumerate}
\item 第~2~章：主定理、遞迴樹、從程式寫遞迴式。
\item 第~7、8、6~章：走訪、由走訪序列重建、建堆與堆積追蹤、中序轉後序。
\item 第~14~章：快速排序的分割與性質、各種排序的比較表（穩定、原地、最壞情況）。
\item 第~12、10~章：開放定址的追蹤、AVL 和紅黑樹的插入。
\item 第~18~章：至少把設計步驟和背包、LCS、矩陣鏈乘練熟。
\item 第~21、24、25~章：BFS／DFS、Kruskal／Prim、Dijkstra 和 Floyd--Warshall 的手算。
\item 第~28~章：NP、NP 完備、歸約方向的觀念題。
\end{enumerate}

每章先讀「本章常見錯誤」和「歷屆試題」，再回頭補不會的節。

\subsection*{路徑三：他系、資管的「資料結構」}

適合：電機、電信、工科、資管等系所單考資料結構。 這類考卷偏資料結構本身，演算法設計和理論很少。

\begin{enumerate}
\item 第~1~章（漸近符號）、第~4~章（陣列位址計算、稀疏矩陣）。
\item 第~5、6、7、8、9、10~章：這是他系考卷的主體，每一章的出現率都超過四成。 第~8~章的二項堆積、左傾樹、費氏堆積在部分學校考得很多。
\item 第~12、14~章。
\item 第~21~章，接第~23~章的割點與雙連通元件、第~24~章、第~22~章的 AOE。
\item 第~11~章略讀（B 樹、trie）。
\end{enumerate}

演算法章在他系考卷裡出現率都不高：第~18~章動態規劃約四分之一，第~16、17、19、20~章和第~26~章以後都在兩成以下，可以最後再看。

\subsection*{路徑四：偏演算法與理論的考卷}

適合：目標考卷常出證明、設計、平攤分析的學校，或想要演算法基礎打得比較深。

在路徑一之外，加重：第~3~章（位能法）、第~17~章（分治的正確性與遞迴式）、第~18、19~章的正確性證明、第~25~章（負邊與 Bellman--Ford）、第~26~章（最大流最小割）、第~28~章（歸約的寫法）、第~16、29~章。 每章的 whybox 都要讀。

\subsection*{路徑五：配合大學課程，從頭讀}

照章的順序讀。 第一次讀可以跳過 whybox 和「實作練習」，期末或準備考試時再回來。 本章的分級表可以拿來決定哪幾章要讀第二遍。

\section*{怎麼看這些數字}
\addcontentsline{toc}{section}{怎麼看這些數字}

\begin{itemize}
\item 數字反映的是過去的出題。 出題老師換人，比重就會變；「少考」不是「不考」。
\item 一題對到哪一節是人工判斷，會有誤差；百分比看大小就好，差一兩個百分點沒有意義。
\item 舊卷（大約 100 學年以前）比較多計算幾何、索引檔、B 樹細節；近年比較多複選題、平攤分析、網路流。 準備時以近 10 年為準。
\item 每章最後的「歷屆試題」收了幾題原題和完整詳解，題目照原卷。 依法令舉行的考試試題依著作權法第~9~條不受著作權保護。
\end{itemize}
